% LINEAR-SWITCH-CUT1
\section*{LINEAR-SWITCH-CUT1: passive switch cuts at linear clause density}

\textbf{STATUS: PROVED/EXACT model theorem plus COMPUTATIONAL random-instance estimates. P versus NP remains OPEN.}

\paragraph{General promise-cut theorem.}
Let \(N=8\binom n3\), and let the input promise be all exact \(m\)-subsets of the \(N\)
possible genuine three-clauses. Consider a passive monotone two-terminal contact scheme whose
connectivity represents a chosen positive class \(P\). Let \(Z\) denote the number of negative
promised inputs.

If an \(s\)-\(t\) cut uses only \(d\) distinct clause labels, then every \(m\)-subset avoiding
those labels disables the cut. Hence
\[
\binom{N-d}{m}\le Z.
\]
Define
\[
\tau_Z=\min\left\{d:\binom{N-d}{m}\le Z\right\}.
\]
When every clause label belongs to a negative promised input, every conducting path must contain
at least two controlled contacts. Edge-Menger therefore gives
\[
C_{\rm contact}\ge2\tau_Z.
\]

For
\[
8\le m\le 7\binom n3,
\]
every clause occurs in both a SAT and an UNSAT promised formula. Thus the two-contact lemma
holds whether connectivity represents SAT or UNSAT.

If a free final output inversion is allowed, the implementation may choose either polarity.
Writing \(S,U\) for the SAT and UNSAT counts,
\[
\boxed{
C_{\rm contact}\ge
2\min(\tau_U,\tau_S)
=
2\tau_{\max(S,U)}.
}
\]
This is the polarity-robust passive cut/Menger bound.

\paragraph{Linear-density asymptotics.}
Let \(m=\alpha n\) and
\[
N=\frac43n^3+O(n^2).
\]
If the majority-class fraction
\[
q=\frac{\max(S,U)}{\binom Nm}
\]
converges to a fixed value in \((1/2,1)\), then for \(d=o(N)\)
\[
\frac{\binom{N-d}{m}}{\binom Nm}
=
\exp\left(-(1+o(1))\frac{md}{N}\right),
\]
and consequently
\[
\tau_{\max(S,U)}
\sim
\frac{N}{m}\log\frac1q
\sim
\frac{4}{3\alpha}\log\frac1q\,n^2.
\]
At exact balance \(q=1/2\),
\[
\boxed{
C_{\rm contact}
\ge
\left(\frac{8\log2}{3\alpha}+o(1)\right)n^2.
}
\]
Thus the count-only cut method has only quadratic scale at a balanced linear-density transition.
If one answer class becomes overwhelmingly dominant and output polarity can be inverted, this
lower bound collapses.

\paragraph{Experiment.}
Uniform simple three-CNFs with no repeated clauses were generated at
\[
n\in\{64,96,128\},
\qquad
\alpha\in\{3,4,4.25,4.5,5,5.25,6,8\},
\]
with 100 formulas per point and deterministic seed 20261003. Each formula was classified exactly
by the project's DPLL solver.

The source and analyzer are
\texttt{experiments/universal-dag/linear\_density\_switch\_cut\_sampler.cpp}
and
\texttt{experiments/universal-dag/linear\_density\_switch\_cut\_analyze.py}.
The raw sample panel is
\texttt{results/raw/universal-dag/2026-10-03\_LINEAR\_SWITCH\_CUT1.csv}.

Two-sided exact Clopper--Pearson 95\% intervals were used only as statistical uncertainty
quantification. They are not finite-family mathematical lower bounds.

At \(n=128\):
\[
\begin{array}{c|c|r|r|r}
\alpha & \widehat p_{\rm UNSAT}
&\text{point switch LB}
&\text{95\% statistical switch LB}
&\text{median DPLL nodes}\\
\hline
3.00&0.00&0&0&32\\
4.00&0.06&660&242&823\\
4.25&0.47&6372&4628&5699\\
4.50&0.86&1430&778&5511\\
5.00&1.00&0&0&2766\\
5.25&1.00&0&0&2366\\
6.00&1.00&0&0&980\\
8.00&1.00&0&0&290
\end{array}
\]

The point lower bound peaks near the empirical SAT/UNSAT crossover. At \(\alpha=4.25\),
\[
6372/128^2\approx0.389,
\]
consistent with the balanced asymptotic coefficient
\[
\frac{8\log2}{3(4.25)}\approx0.435.
\]

\paragraph{Cross-model hostile controls.}
The existing min-fill panel at \(n=128\) grows smoothly from width about 76 at \(\alpha=3\)
to 91 at \(\alpha=5\), rather than reproducing the sharp switch-cut peak.
The exact width-three resolution panel has no width-three refutations among five tested
\(n=128\) instances at each of \(\alpha=5.25,6,8\), even though the polarity-robust passive
cut bound is zero on the 100-sample panel there.

Therefore the scalar count-only passive cut resource is \textbf{FALSIFIED/RETIRED as a candidate
common hardness invariant}. It mainly measures class imbalance. It cannot explain the
resolution-hard UNSAT regime and has only quadratic scale near balance.

\paragraph{Literature context.}
Friedgut's work establishes sharp-threshold behavior for random \(k\)-SAT in the appropriate
threshold-sequence sense. For small fixed \(k\), including \(k=3\), the exact limiting
satisfiability threshold remains unresolved. Rigorous bounds for random 3-SAT remain separated,
while statistical-physics and empirical work place the transition near \(4.26\).
Recent ECCC work gives an \(O(n/\log n)\) scaling-window upper bound for random \(k\)-SAT for
every fixed \(k\ge3\).

Jukna's contact-circuit entropy method is directly relevant to the required repair: a useful
cut resource must capture inner information/local dependence, not only the cardinality of the
two answer classes.

\paragraph{Next target.}
Replace the scalar negative-count \(Z\) with a structural cut matrix
\[
M_{h,z}=F(h\cup z),
\]
where rows are multisign partial states on one side of the cut and columns are legal
continuations on the other. Test:
\[
\operatorname{rank}_{\mathbb F_2},
\quad
\text{rectangle/fooling-set bounds},
\quad
\text{communication complexity},
\quad
\text{multisign/projective dimension}.
\]
The target is a cut lower bound that survives severe SAT/UNSAT imbalance and therefore measures
nonmergeability rather than class balance.

\begin{thebibliography}{9}
\bibitem{Jukna88}
S.~P.~Jukna,
\emph{Entropy of Contact Circuits and Lower Bounds on Their Complexity},
Theoretical Computer Science 57 (1988), 113--129.

\bibitem{Friedgut99}
E.~Friedgut,
\emph{Sharp Thresholds of Graph Properties, and the \(k\)-SAT Problem},
Journal of the American Mathematical Society 12 (1999), 1017--1054.

\bibitem{HajiaghayiSorkin}
M.~Hajiaghayi and G.~Sorkin,
\emph{The Satisfiability Threshold of Random 3-SAT Is at Least 3.52},
IBM Research Report RC22942, 2003.

\bibitem{DuboisEtAl}
O.~Dubois, Y.~Boufkhad, and J.~Mandler,
\emph{Typical Random 3-SAT Formulae and the Satisfiability Threshold},
Theoretical Computer Science, upper-bound work yielding 4.506.

\bibitem{Carenini26}
G.~Carenini,
\emph{The Scaling Window of Random \(k\)-SAT},
ECCC TR26-145, 2026.
\end{thebibliography}

% END_OF_NOTE::LINEAR-SWITCH-CUT1
